\section{No intervention has passed the gate}
\label{sec:causal-results}

\subsection{Where the tool choice is carried}
\label{sec:map}
Before steering, we mapped where the choice of \texttt{answer} is made. On two \fsp{} and two \fsm{} completion steps, we patched the residual stream from the same task's previous step (whose recorded action is a click) at each of \cfgBigLayers{} layers and four positions, \cmProbes{} patches in all, with the original reasoning held fixed. The quantity is the logit margin between the first token of \texttt{answer} and that of the best other tool. Identity patches give exactly zero.

The decision is invisible at the end of the prefill and barely visible mid-reasoning: the largest mean margin change is \cmPrefillPos{} and \cmMidPos{} logits there, for \fsp. It appears only at the tool-name position, from about layer 16 upward, where patching moves the margin by up to \cmMinPos{} logits on \fsp{} and \cmMinNeg{} on \fsm{} (Fig.~\ref{fig:map}). Component scans single out three attention heads, L19/H15, L23/H1 and L31/H14, and ten late MLP neurons. Their attention falls mostly on the system-prompt definition of \texttt{answer} (L19/H15, L23/H1) and on the JSON being written (L31/H14). % source: data/pod_astra/circuit/REPORT.md "Où regardent les têtes ?"
Removing all of them together lowers the margin by \cjDeltaMin--\cjDeltaMax{} logits, and flips \cjFlips{} of the \cjN{} decisions (smallest remaining margin \cjMarginMin). Across all \cmCompProbes{} component patches and \cjAll{} joint patches, \cmCompFlips{} and \cjAllFlips{} decisions flip: the components move the margin, not the action. The only flips come from residual patches (\cmFlipsPos{} on \fsp{} and \cmFlipsNeg{} on \fsm{} of \cmProbes), at the tool name in the last layers, where replacing the residual stream transplants the donor's choice into the output head. They show where the choice is carried, not where a false judgement is formed. % source: data/pod_astra/circuit/REPORT.md "Carte causale"
 The effect is the same on true and false completions (Fig.~\ref{fig:map}c). It is a generic tool-retrieval circuit, not a false-success circuit.

\paragraph{A ceiling test.} The patches above replace one layer at a time; a ceiling test asks how much must be replaced at once. With the recorded input forced and only logits read, the outputs of whole blocks of writers are swapped in at the decision position from a donor step where the model clicked: the same task's previous step for two \fsp{} contexts, another task's step for one true completion (Fig.~\ref{fig:ceiling}). On the first \fsp{} context (margin \ceilBase), all 32 MLPs together bring the margin to \ceilMlpAll{} without flipping it, while the MLPs and token mixers (attention and DeltaNet) of layers 12--23 flip it (\ceilBothMiddle), as do those of layers 24--31 (\ceilBothLate), and those of layers 0--11 do not (\ceilBothEarly). The second \fsp{} context behaves the same (MLPs \ceilSecMlpAll, block 12--23 \ceilSecBothMiddle), and so does the true completion (block 12--23: from \ceilTrueBase{} to \ceilTrueBothMiddle). Swapping with itself and swapping the final residual reproduce the baseline and the donor exactly. % source: data/pod_astra/broad/ceiling.jsonl, ceiling_resume.jsonl, ceiling_checks.json (asserted in figs/broad.py)
Three caveats bound this result. The first \fsp{} context is an infeasible task, on which \texttt{answer} (with a refusal) is the correct tool (Appendix~\ref{app:infeasible}). The donor differs in screen, reasoning, thought and position, not only in decision. And since a true completion flips too, the test locates where the identity of the next tool token is assembled, not a mechanism of false success. No particular block is needed: the MLPs and mixers of layers 12--23 or of layers 24--31 each flip the first context, and removing a single direction of the swap difference from the writers of layers 12--23 brings the margins of the two \fsp{} contexts to \projOneFspA{} and \projOneFspB{}, but also those of two true completions to \projOneTrueA{} and \projOneTrueB. % source: data/pod_astra/broad/projection_summary.json (rank 1, dose 1)
 It is consistent with the failure of every edit confined to one or two layers or to the MLPs, and it says nothing about which edit would separate false from true completions.

\paragraph{The tool call follows the text already written.} An anchoring test keeps the screen, the history and the weights, and replaces only the text the model has generated at the decision step with the donor's. On the first \fsp{} context (margin \ancAOrig), the donor's reasoning alone lowers the margin to \ancAReason, the donor's JSON prefix alone (note and thought) to \ancAJson, and both together flip it (\ancABoth). The second context behaves the same (\ancBOrig; reasoning \ancBReason, JSON prefix \ancBJson, both \ancBBoth). % source: data/pod_astra/broad/anchoring.jsonl, anchoring_summary.json, anchoring_plan.json
Three limits weigh on this. There is no control on true completions or with a neutral text. Only the second context is feasible, and the donor is the previous step, whose text announced a click, so that copying its text largely copies its decision. And the text does not account for everything: with all of it swapped, the margin stays \ancAScreen{} and \ancBScreen{} logits above the donor's, a share that comes from the screen and the history. Within these limits, the tool call depends causally on the text already written, the note and thought weighing most, and the \margPosMin--\margPosMax-logit margin of \S\ref{par:ceiling} is largely conditioned on that text rather than an intrinsic robustness of the decision.


\placefig{fig-ceiling}

\placefig{fig-map}

\subsection{Every steering arm}\label{sec:every-arm}
Table~\ref{tab:causal} lists every intervention run on the 4B false-success target, each under the gate of \S\ref{sec:protocol} and on the same six screening contexts. Each line corresponds to a distinct hypothesis taken from the literature.
\begin{itemize}
\item \emph{Tool choice is linearly steerable} \citep{wu2026toolcalling}: add the position-matched \fsp{}$-$\fsm{} direction at the last prefill token, at layer 16 ($\alpha=\cAlphaSixteen$ times its norm) and layer 24 ($\alpha=\cAlphaTwentyFour$).
\item \emph{The decision is taken during reasoning} (\S\ref{sec:think}): add a ``conclusion'' direction, fitted on the reasoning tokens, at every reasoning token of layer 31.
\item \emph{False success is excess expected value} \citep{jiang2026value,kumaran2026commit}: push toward low value or low commitment at layer 24 or 18.
\item \emph{Late MLP neurons write the decision} \citep{dutta2026termination}: zero the top-16 or top-64 neurons by direct attribution to the \texttt{answer} logit.
\item \emph{The model fails to check}: add a ``verify'' versus ``conclude'' direction (contrastive activation addition) at layer \verifyLayer{} at 1 and 2 times the residual norm, over the prefill and the first 20 generated tokens.
\item \emph{False success is a perception of completion}: subtract the imminent-completion direction of \S\ref{sec:perceive} at layer \pcLayer.
\item \emph{The model opens with its conclusion} (\S\ref{sec:think}): raise the output row of the token ``Wait'' for the states that open a completion, a \waitWeights-weight edit of the output head.
\item \emph{Memory of the early episode fades}: lengthen or shorten the retention of the nine late linear-attention layers by a factor 1.25.
\item \emph{The circuit of \S\ref{sec:map} carries it}: scale its three heads and ten neurons (C1) by $1-\alpha$ with $\alpha=0.5$ and $1$, against a random direction of the same contribution norm.
\end{itemize}
No arm passes the gate. Most leave the false-``done'' rate at its original value, including removing the whole C1 circuit ($\alpha=1$), which answers on every false completion, like the random control. The neuron arms, screened on six other disjoint contexts whose own baseline is \cNeuronBase{} \fsp{} answers, were stopped by a mathematical bound: with \cNeuronObs{} answers already observed on the first unseen \fsp{} task, at most \cNeuronMax{} answers out of 8 would have been admissible. The value-axis arm did change the text. An independent reader found an explicit intention to re-evaluate in \doubtPushFsPos{} \fsp{} reasoning traces against \doubtBaseFsPos{} at baseline, but \doubtPushStrict{} verification actions and the same answer rate. % source: data/pod_astra/value/*_VERIFICATION_ANNOTATIONS.json (secondary/strict flags)
The model wrote that it should check, then answered anyway.

The ``Wait'' row is the one edit that lowers false answers without touching anything else: \waitFsp{} on false completions, with \waitFsm{} true completions kept, \waitHealthy{} healthy actions valid and \waitText{} general-text controls token-identical. It still fails the 50\% relative criterion, and an independent review rates none of the two replacement actions useful: one reloads a settings page for an invented reason, the other scrolls without fixing the dates. % source: data/pod_astra/mit1/REPORT.md "Lot6 wait2"; data/pod_astra/mit1/semantic_lot06.json
Making the model pause is not enough to make it repair.
A conditional variant adds an early end-of-reasoning row to the same edit (\earlyExitWeights{} weights). Inside its construction set, \hfourInFalse{} false completions still answer. On one Budget sample the model reasons correctly that next month is December, says the dates must be changed, and generates a click on the date summary of the results page; it was not executed. % source: data/pod_astra/mit1/semantic_lot14.json; data/pod_astra/mit1/REPORT.md first section
These are intentions to correct, not corrections: an independent review of the screenshot finds that the dates sit in plain summary text, next to a separate ``Modify Rental Details'' control, and that no step proves the summary opens an editor. All \hfourIntentions{} such generated clicks, over \hfourIntentLots{} lots of this edit and its diagnostics, are reclassified as unverified, \hfourConfirmed{} confirmed useful; no effect on the interface or on the OSWorld score is shown. % source: data/pod_astra/mit1/review_budget_affordance.{json,md} (sha 0aa1bc6a); data/local/JOURNAL.md entries 00:15 and 00:20 UTC
Nor does the edit transfer. On nine training tasks left out of its construction, the final audit counts \hfourOutUseful{} useful actions on false completions (original \hfourBaseUseful), \hfourOutErrId{} outputs that name the real error (original \hfourBaseErrId), and \hfourOutFalse{} false finishes (original \hfourBaseFalse), while true completions go from \hfourBaseTrue{} to \hfourOutTrue{}. % source: data/pod_astra/mit1/semantic_lot19.json (sha 44a9011b, supersedes the first audit's 1/12 useful and 5/12 identified); SEMANTIC_AUDIT_LOT15_ADDED_BASELINE.json
The edit leaves a mechanistic signal, the model recognising the error then rationalising it away, and no demonstrated useful effect.

\paragraph{These screens are pilots.} Each efficacy cell holds 8 generations from two tasks. 8/8 has a Wilson 95\% interval of [\wilsonEightLo, 100]\%, and samples clustered in two tasks make the real uncertainty larger, so a rejection means only that no large effect appears on these contexts, not that none exists. By design the gate flags a reduction to at most \gateMaxAnswers/8 answers; smaller effects pass unnoticed. Table~\ref{tab:causal} gives the interval of every efficacy cell.

\paragraph{A ceiling effect is likely.}\label{par:ceiling} At the tool-name position the margin of \texttt{answer} over the best other tool is \margPosMin--\margPosMax{} logits on false completions and \margNegMin--\margNegMax{} on true ones. This position comes after the JSON \texttt{thought} field, which the output schema writes before the tool call and which already states that the task is complete, so the margin measures how confidently the model copies a decision it has already verbalised, not the margin of the decision itself. % source: ~/osw/rtg/Agent/holo_agent/tools.py make_step_model (note, thought, tool_call); data/local/reports/REDTEAM_0305.md point 3
Even so, the perturbations tested are small: next to the residual stream, they are \doseSixteen{} (layer 16) and \doseTwentyFour{} (layer 24) of the mean residual norm for the matched direction, \doseThink{} for the reasoning-token direction (difference-of-means norm \thinkScale{} against a residual norm of \thinkRThirtyOne{} at layer 31), and 0.5 for the value axes, % source: data/pod_astra/value/REPORT.md "Norme = norme résiduelle moyenne"
while removing the whole C1 circuit lowers the margin by \cjDeltaMin--\cjDeltaMax{} logits without flipping it. Larger doses do not break the ceiling either. The ``verify'' direction at one and two times the residual norm leaves \fsp{} at 8/8; at dose 2 a healthy context loops on ``I see'' up to the 3{,}072-token cap, which stopped the arm after \verifyRejCompleted{} of 24 outputs. % source: data/pod_astra/verify/verify_a2.rejected.json; JOURNAL.md 22:10 UTC astra-verify
Dose sweeps at layers 24 and 31 with random controls were not run.

A second target, getting stuck at the step limit, shows the opposite failure of steering. A residual addition at layer 16 or 24 damaged healthy actions before any stuck context was generated; after \stuckSixteenObs{} outputs, the best reachable health rate was \stuckSixteenBest\%, below the required \stuckSixteenReq\%. % source: data/pod_astra/stuck/CAMPAIGN_STATUS.json

\placefig{tab-causal}

\subsection{A detector of imminent completion beats position, not false success}
\label{sec:perceive}
All directions above contrast false and true completions, which differ in position. \texttt{astra-perceive} sidesteps this by learning a direction \emph{inside successful episodes}: the last-prefill residual of the final, correct \texttt{answer} step minus that of the step just before it, on \pcPairs{} pairs from \pcTasks{} construction tasks (reserve excluded, cross-validation grouped by task across both passes). Because both members of a pair sit next to each other in the same episode, position cannot separate them across tasks. The direction peaks at layer \pcLayer{} with task-grouped AUROC \pcAuc, against \pcBaseStep{} for the step index and \pcBaseCtx{} for context length (Fig.~\ref{fig:perceive}a). This is the first signal in this study that position does not explain, in contrast to the false-success direction of Table~\ref{tab:fs-auc}. Within a pair, both baselines are trivially perfect, so the pooled AUROC is the relevant comparison, and the layer was chosen on the same cross-validation.

Read on \pcMatchN{} false/true completion pairs matched on pass, step and context length, with each task's direction fitted without that task, the direction rises at the ``done'' step in \pcFsPosPositive/\pcMatchN{} false and \pcFsNegPositive/\pcMatchN{} true completions (mean \pcFsPosMargin{} and \pcFsNegMargin{} standard deviations; Fig.~\ref{fig:perceive}b). At the terminal step, false completions sit \pcTermDiff{} standard deviations from true ones, above them in only \pcFsPosGt/\pcMatchN{} pairs. It is therefore a detector of imminent completion, not of false success: it fires before a false ``done'' as it does before a true one. On these \pcMatchN{} pairs, the projection reads a false completion as complete. This is compatible with a perception error: the screen is taken as showing a finished task. Two same-task image pairs from construction data illustrate it, a spreadsheet cell claimed green whose text has no green pixels, and rental dates claimed for next month while the page shows the current one. % source: data/pod_astra/perceive/REPORT.md "Screen and visual evidence"
It does not separate misperception from commitment to the action. Subtracting it at layer \pcLayer{} (up to \perceiveMaxDose{} of the residual norm, on its own screening contexts whose reference is \perceiveFsBase/8 \fsp{} answers) gives \perceiveFsSeries{} \fsp{} answers out of 8 for $\alpha=0.5,1,2,4$, against \perceiveFsRandom/8 for a random direction of the same norm, and no dose passes the gate (Table~\ref{tab:causal}). The few drops are not repairs: the non-answer actions that replace \texttt{answer} are empty or invented text entries, such as a made-up e-mail address. % source: data/pod_astra/perceive/REPORT.md "Génération causale"
Attention maps on two same-task pairs found no head, out of \perceiveHeadsTested{} tested, that reads the evidence region more when the model succeeds (\perceiveHeadsEligible/\perceiveHeadsTested{} pass the frozen rule). The perception hypothesis is not refuted by this small screen, but it has no support beyond the projection of Fig.~\ref{fig:perceive}.

\placefig{fig-perceive}

\subsection{Edits without a direction: extrapolation and the Qwen post-training vector}
\label{sec:expo}
Direction-free weight edits do not depend on locating anything. Holo-3.1-4B descends from the post-trained Qwen3.5-4B, not from its base: in the language layers, the fine-tune moves the weights by \linHoloQ\% of their norm, against \linHoloBase\% from the base and \linQBase\% between Qwen's base and post-trained models, and H Company's delta is almost orthogonal to Qwen's post-training (cosine \linCos). % source: data/local/expo/EXPO_REPORT.md section 2
We therefore tested model extrapolation, $W(\lambda)=W_{\text{Qwen}}+\lambda\,(W_{\text{Holo}}-W_{\text{Qwen}})$ \citep{zheng2024expo}, on all language layers, attention only or MLPs only, and the addition of half of Qwen's post-training delta, in the spirit of reasoning vectors \citep{zbeeb2025reasoningvectors}. None changes false answers: they stay at 16/16 in every arm but one (\expoFspMin/16, MLPs only), where true completions drop more (Table~\ref{tab:expo}). At $\lambda=1.5$ one output in 32 loops to the token cap. The edit at $\lambda=1.25$ keeps text and JSON intact, and it runs online as an arm of Table~\ref{tab:b4-heldout}, because only OSWorld measures whether actions improve.

\placefig{tab-expo}

\subsection{Mechanism-level edits change the text, not the tool call}
\label{sec:attractor}
Three more gradient-free edits target a mechanism rather than a direction. H2 resets the 48 memory-gate matrices of the 24 Gated DeltaNet layers to Qwen's; H3 moves the 96 language MLP matrices halfway to Qwen's, leaving attention, memory, vision and the output head untouched; H5 moves the input embedding of \texttt{</think>} halfway toward the double-newline row, with the output head untied and kept. % source: data/pod_astra/free/REPORT.md sections H2, H3, H4 and H5
Each was served by vLLM on separate hardware against a re-served original, with the same 8 seeds on the same 6 development contexts. Each reproduces the original's tool counts exactly: \attHtwoFsp{} false and \attHtwoFsm{} true completions answer, every healthy action stays valid, no output is invalid (Table~\ref{tab:attractor}). Only \attHtwoFull, \attHthreehalfFull{} and \attHfivehalfFull{} outputs are token-identical to their paired original, and \attHtwoReason, \attHthreehalfReason{} and \attHfivehalfReason{} up to the end of reasoning (Fig.~\ref{fig:attractor}). This does not by itself show that the edits reach the computation: we did not measure how often two runs of the unmodified model with the same seeds differ on a batched server, and the H5 edit acts only once \texttt{</think>} has been generated, so its identical reasoning prefixes are expected. What the table does show is that the tool call is robust to these edits. % source: data/local/reports/REDTEAM_0305.md point 3; data/pod_astra/free/REPORT.md (batched execution need not be identical across runs)

One edit does move the tool call, and shows why a moved tool is not a correction. Scaling the attention scores of the eight full-attention layers by $s=0.5$ (H6, a softer attention, 8 \texttt{q\_norm} vectors) lowers false-completion answers to \attHsixFsp, while true completions answer \attHsixFsm{} and healthy actions stay valid. All \infHsixNonAnswer{} non-\texttt{answer} calls fall on the infeasible context, where \texttt{answer} is the correct tool; on the feasible \fsp{} context H6 answers \infFreeHsixExcl, like the original (Appendix~\ref{app:infeasible}). % source: data/pod_astra/free/REPORT.md "H6: full-attention temperature"; data/pod_astra/free/external/free_h6_s05/summary.json
A blind semantic audit, with original and edited outputs mixed, finds \attHsixAuditUseful{} useful actions, as for the original, and \attHsixAuditFalse{} outputs that still claim success (original \attOrigAuditFalse); the \attHsixTheme{} non-\texttt{answer} calls are unhelpful theme clicks (GTK or Classic) justified by an incorrect prediction that they will restore the old interface, and none is useful. The tool moved; the task did not get closer to done. % source: data/pod_astra/free/audit_external_h6_pair.fs_text.adjudicated.verdict.json and .adjudication.md (aggregate.json predates the adjudication); data/local/JOURNAL.md entry 01:50 UTC
The second pre-registered dose, $s=0.25$, was stopped at its pilot: the only output runs to the \hsixQuarterTokens-token cap without closing its JSON, which closes the family. % source: data/pod_astra/free/external/free_h6_s025/pilot.json; data/pod_astra/free/VERIFY_H6_QUARTER_ABORT.md
Four last edits close the series. Doubling the H2 gate edit ($\lambda=2$, past Qwen), moving all 401 language tensors 20\% of the way to Qwen at once (H9, a coordinated restoration rather than one block at a time), and halving the share of the query gain carried by the 64 rotary coordinates of each full-attention layer, the norm of the gain vector kept (H10, $s=0.5$), leave the decision in place: false completions answer \attHtwotwoFsp, \attHnineFsp{} and \attHtenhalfFsp, true completions \attHtwotwoFsm, \attHnineFsm{} and \attHtenhalfFsm, and the blind audit finds no useful action. Removing that share entirely (H10, $s=0$) repeats H6. False-completion answers fall to \attHtenzeroFsp, all on the infeasible context, where \attHtenzeroInfClick{} theme clicks and \attHtenzeroInfWait{} wait replace an answer that should have explained the impossibility; the feasible context still answers \infFreeHtenzeroExcl, and \attHtenzeroHealthyTrunc{} healthy steps end in truncated outputs, a loss of at least \attHtenzeroHealthyLoss{} points that rejects the edit. % source: data/pod_astra/free/REPORT_final_0505.md (H2, H9, H10), FINAL_RESULTS_TABLE.md; external/free_{h2_l2,h9_a020,h10_half,h10_zero}/summary.json; audit_external_{h2_l2,h9,h10_half,h10_zero}_pair.aggregate.json; audit_external_h10_zero_pair.fs_text.verdict.json; H10_ZERO_EARLY_REJECTION.json

Merging in a donor trained on desktop tasks moves the decision globally, not selectively. Adding a fraction $\beta$ of the task vector of CUA-Lite, a Qwen3.5-4B fine-tune on ScaleCUA desktop data, to the language layers of Holo ($W_{\text{Holo}}+\beta\,(W_{\text{CUA}}-W_{\text{Qwen}})$), and serving it with the same 51 calls as the exact original model, re-served on the same GPU type and vLLM version (in a separate account and container), lowers false-completion answers to \mergeCuaQFsp{} at $\beta=0.25$ and \mergeCuaFsp{} at $\beta=0.5$, and true-completion answers further, to \mergeCuaQFsm{} and \mergeCuaFsm{} (original \mergeCuaOrigFsp{} and \mergeCuaOrigFsm), with every healthy action valid and no invalid output; on the feasible \fsp{} context alone, $\beta=0.5$ answers \infMergeCuaHalfExcl. The dose--response is a global drop of \texttt{answer}, not a discrimination. An independent audit finds the replacing actions coherent and unnecessary: scrolls after the model has recognised success, and verification clicks after true completions. The model continues instead of concluding, whatever the truth of the completion. A donor trained on filtered web tasks, Fara1.5-4B, leaves the decision in place at $\beta=0.5$ (\mergeFaraFsp{} and \mergeFaraFsm). BF16 rounding blurs small merges: at $\beta=0.5$, \mergeCuaUnchanged{} of the \mergeCuaLangTensors{} language tensors are left unchanged and the update norm is \mergeCuaActualNorm{} against an intended \mergeCuaIntendedNorm; at $\beta=0.25$, \mergeCuaQUnchanged{} tensors are unchanged. CUA-Lite's training data may overlap OSWorld tasks. % source: data/pod_astra/merge/STATIC_RESULTS.json; VERIFY_MODAL_REFERENCE.{md,json} (original revision, H100, vLLM 0.23); VERIFY_CUA_B050_BEHAVIOR.{json,md}, VERIFY_CUA_B025_BEHAVIOR.md, VERIFY_FARA_B050_BEHAVIOR.md; REPORT.md (donor origin, contamination limitation)
Online, the merge shows no detectable gain over an original served by the same endpoint: \onAllCuaDelta{} points over \onAllCuaN{} tasks (95\% CI \onAllCuaCi; \S\ref{sec:online}).

A closed-form edit of the writers that the ceiling test points to moves the decision as well, and just as unselectively. In the spirit of AlphaEdit \citep{fang2024alphaedit}, the \alphaNWriters{} residual writers of layers \alphaLayers{} (MLP down-projections and mixer output projections) receive a ridge update that maps \fsp{} keys (constructed on the ceiling-test contexts, the first of them infeasible) toward the donor displacement while being projected off the span of protected keys (true completions, healthy steps and text prompts). On contexts outside its construction, served next to the original on the same GPU, false-completion answers fall from \alphaOrigFsp{} to \alphaEditFsp{} but true-completion answers fall further, from \alphaOrigFsm{} to \alphaEditFsm, with \alphaEditHealthy{} healthy actions valid and \alphaEditInv{} invalid outputs. The whole \fsp{} drop comes from one of the two contexts (\alphaOrigFspCtxB{} to \alphaEditFspCtxB; the other stays at \alphaEditFspCtxA). A first visual audit of the replacing actions, which had described the screenshot of the previous step, was retracted; the corrected audit finds that on the e-mail true completion all four recipients were already complete, so the three generated \texttt{write} continuations there would potentially damage a correct state (they were not executed). % source: data/pod_astra/broad/{eval_original,eval_p100}/summary.json, derive_residual.json; data/local/JOURNAL.md 02:14, 02:35 and 02:37 UTC (semantic_p100.{json,md} is the retracted audit)
A random update on the same writers, with the same right factors and its per-writer BF16 perturbation norm matched to the learned one within 0.1\%, leaves everything in place (\fsp{} \alphaRandFsp, true completions \alphaRandFsm, healthy \alphaRandHealthy): the effect is specific to the learned direction, a real causal lever on stopping, without selectivity. % source: data/pod_astra/broad/eval_random100/summary.json; data/pod_astra/broad/REPORT.md (random controls)
At half the dose the edit gives \alphaHalfFsp{} and \alphaHalfFsm: at both doses, true completions fall as fast as or faster than false ones. % source: data/pod_astra/broad/eval_p050/summary.json
A 25\% dose was not run.

Two independent methods, a merge and a closed-form edit, thus move the decision, and in both cases true completions fall at least as much as false ones, on two true-completion contexts: each finds a lever without a discriminator. A simple argument makes this unsurprising. An edit applied at or after layer $\ell$ acts on a function of the residual stream at layer $\ell$ (all positions, and the recurrent state of the linear-attention layers), so it can treat a false and a true completion differently only to the extent that some function of that stream separates them, and to be useful that function must generalise to tasks it was not built on. Our probes, which read the stream at two positions, found no such separation beyond episode length (Table~\ref{tab:fs-auc}).
A direct test on the writers that both methods edit finds the same. On \rmFsp{} false and \rmFsm{} true completions outside the reserve and the evaluation contexts, infeasible tasks excluded, a ridge probe on the inputs (keys) of the \rmWriters{} writers of layers 12--23, cross-validated by task, does no better than position. Averaged within test folds and weighted by their \rmPairs{} positive--negative pairs, its median AUROC is \rmPreAuc{} (range \rmPreRange) at the end of the prefill and \rmEndAuc{} at the end of reasoning, against \rmBaseStep{} for the same probe on the step index; pooled out of fold, it is \rmPreOof{} and \rmEndOof{} against \rmBaseStepOof. % source: data/pod_astra/regmean/overlap.json (pair_weighted_fold_auc and oof_auc, cv_probe in astra-regmean/analyze.py), overlap_summary.json; manifest.json (sha f517da46)
The two key subspaces are not identical (median principal angle \rmPreAngle--\rmEndAngle$^\circ$), but \rmEndEnergy--\rmPreEnergy\% of the false-completion key energy lies in the true-completion subspace. This is an empirical, non-total overlap together with a strong confound of position and length, not an impossibility: a low AUROC for a linear probe is not an information bound, the probes do not exhaust the stream, and even identical subspaces would not force a selective edit or merge to act globally, which would require proportional covariances. % source: data/pod_astra/regmean/overlap_summary.json, overlap.json (its own caveat)


\placefig{fig-attractor}

The same picture comes from three other directions. Seven two-neuron conditional circuits (\mitTwoWeights{} weights each), written into the MLP of layer 19 or 31 to fire on an imminent conclusion and push toward verification, give \mitTwoFsp{} false answers; the layer-31 version inserts one ``Wait'' in every output, after which the model sometimes recognises the wrong month on Budget and rationalises it. % source: data/pod_astra/mit2/REPORT.md
Even the input such a circuit aims for does not help: injecting a generic instruction to compare each requested detail with the visible evidence into the reasoning, with the original weights, late or early, gives \mitTwoInj{} false answers. A narrow MLP edit adapted from the weight-space programming recipe (64 columns in each of layers 17 and 20, \ratioWeights{} coefficients targeted) gives \ratioSeries{} premature answers at $\alpha=0.5$, 1, 1.5 and 6, and \ratioUsefulAll{} useful actions in any arm; a random direction at $\alpha=6$ gives \ratioRandSix, and true completions stay at \ratioTrueRange{} against \ratioTrueOrig. % source: data/pod_astra/ratio/semantic_summary.json; data/pod_astra/ratio/REPORT.md "Construction figée", "Résultats finaux"
The last hypotheses closed the same way. Promoting the end-of-thinking token so that the model reasons less leaves \texttt{answer} rates unchanged at the dose that keeps outputs valid (false completions \shortTwelveFsp{} against \shortOrigFsp, true completions \shortTwelveFsm{} against \shortOrigFsm, on a 28-context panel), and breaks the JSON at higher doses (\shortSixteenInv{} and \shortTwentyInv{} invalid outputs before the arms were stopped). In \mitOneLateN{} further edits (penalising completion words, amplifying three task-reading heads, interpolating late MLPs toward Qwen), every dose gives 0/8 useful actions. An adaptation of TACT \citep{sui2026tact}, which steers overthinking and overacting at the end of reasoning, could not be built as specified: after a blind audit of the labels, \tactOT{} overthinking and \tactOA{} overacting states remain, and the overthinking axis (AUROC \tactAuc) barely beats position and reasoning length (\tactPosLen). % source: data/pod_astra/short/REPORT.md; data/pod_astra/mit1/REPORT.md H7-H9; data/pod_astra/tact/REPORT.md "First signal and audited freeze v3"
Pushing the MLP edit all the way breaks the model rather than the decision: with the MLPs set exactly to Qwen's, \attHthreefullInv{} outputs are invalid (its only \texttt{answer} calls are on the infeasible context), and one pilot output repeats a single token over \attHthreeRepeatPct\% of its 3{,}072 tokens. % source: data/pod_astra/free/REPORT.md H3full

These screens share the limits of \S\ref{sec:every-arm}: 16 samples per stratum come from two contexts, the circuit arms stopped after 5 outputs once the gate could no longer be met and never measured protected classes, and the MLP edit is an adaptation of the published recipe, not a numerical reproduction. % source: data/pod_astra/mit2/REPORT.md "Résultats et dénominateurs"; data/pod_astra/ratio/REPORT.md "Protocole et sources"

\placefig{tab-attractor}

\subsection{Compiling an intervention into the weights}
\label{sec:compile}
No steering candidate passed the gate, so none was compiled; the weight edits above were built directly by their own constructions. We state the compilation step because it constrains which steering interventions are worth testing. Let the residual stream at layer $\ell$ be $x_\ell = e + \sum_{j<\ell}(a_j + m_j)$, where $a_j=W^O_j z^a_j$ and $m_j=W^{\text{down}}_j z^m_j$ are the attention and MLP writes and $e$ the token or image embedding, and let $P=I-uu^\top$.

\begin{proposition}[Directional ablation equals writer orthogonalisation; cf.\ \citep{arditi2024refusal}]\label{prop:equiv}
In a pre-norm decoder whose sublayers read the residual stream only through their normalisation and whose writers have no bias, replacing every writer $W\in\{W^O_j, W^{\text{down}}_j, W^V\}$ by $PW$, the embedding by $W_EP$ and the vision-merger bias by $Pb^V$ gives a forward pass identical to the original model with $x\leftarrow Px$ applied after every write, at every layer and position.
\end{proposition}
\begin{proof}[Proof sketch]
By induction over writes: if $x$ has no $u$-component, every normalised sublayer input is unchanged, so every $z$ is unchanged, and each new write $PWz=P(Wz)$ keeps the sum orthogonal to $u$.
\end{proof}
Four consequences follow. A hook at block outputs is not equivalent, because the MLP then reads the un-ablated attention output. $W\leftarrow(I-\alpha uu^\top)W$ with $\alpha\neq1$ scales each write's component by $1-\alpha$, whereas a residual hook compounds over layers. Tied embeddings (\cfgBigTied{} in both models) mean that projecting $W_E$ also projects the output head unless it is untied. And an addition $h\leftarrow h+\beta u$, like every intervention in Table~\ref{tab:causal} except projections, can be compiled only as a bias, which the bias-free projections of Holo-3.1 do not have. The circuit arm is the exception: scaling heads and neurons is exactly a column scaling of $W^O$ and $W^{\text{down}}$, 1{,}991{,}680 weights or 0.0439\% of the model, and it has been tested in generation with no effect. % source: data/pod_astra/circuit/REPORT.md "Édition des poids et admission"
Any compiled patch must reproduce the steered behaviour within 5 points per stratum before serving. The serving path is ready: a patched checkpoint is loaded by vLLM in about seven minutes, a null patch changes 0/724 tensors and reproduces the original outputs. % source: JOURNAL.md "Validation en ligne (online-eval) — prête"
Since no candidate passed, no compiled edit or equivalence check was run.

\subsection{Online evaluation}
\label{sec:online}
No candidate passed the gate, but five arms were run online, each paired by task with an original run on the same harness, and the arms served by Modal with an original served by the same Modal endpoint (Table~\ref{tab:b4-heldout}). The original itself is the first measure of noise. On the confirmation reserve it scores \onReserveOrigPct\% on one harness and, on the other, \onReserveGhOrigPct\% when served by the pod and \onReserveModalOrigPct\% when served by Modal: the same weights differ by \onNoiseDelta{} points across harnesses (\onNoiseDisc{} tasks better and worse) and by \onNoiseServerDelta{} points across servers (95\% CI \onNoiseServerCi), the size of every effect below. Against it, no arm shows a gain: ExPO at $\lambda=1.25$ \onReserveExpoDelta{} points (95\% CI \onReserveExpoCi), the \texttt{answer} output-row edit \onReserveLmrowArm{} against \onReserveLmrowOrig, the corrective prompt \onReservePromptArm{} against \onReservePromptOrig, and, against the Modal-served original, the CUA-Lite merge \onReserveCuaDelta{} (\onReserveCuaCi; \onReserveCuaDisc{} tasks better and worse) and the closed-form edit \onReserveBroadDelta{} (\onReserveBroadCi). % source: data/local/online/ (figs/online.py)
Two reserve tasks score 0 under the pod-served original and 1 under both Modal-served arms; they score 0 under the Modal-served original as well, so the gap came from run-to-run variation, not from the server. The harnesses do differ: two \texttt{os} tasks score 0 on the orbit harness and 1 on the GitHub one, so runs from different harnesses are never pooled. The merge also produced 6 unparsable steps against none for the pod-served original, a sign that it degrades the output format. % source: data/local/online/ (figs/online.py, SERVER_SENSITIVE asserted 0 under results_confirm_gh_orig_modal_confirm); data/local/reports/REDTEAM2_0335.md point 2
The decision rule for the merge was fixed before its results were read: the merge and an original served by the same Modal endpoint are compared on the 45 reserve and 120 \textsc{test} tasks, and a gain is claimed only if the lower bound of the task-bootstrap 95\% CI of the paired difference is above zero. By the pre-registered rule (lower 95\% bound above zero), the merge shows no detectable gain: over the \onAllCuaN{} paired tasks it scores \onAllCuaArm{} against \onAllCuaOrig{} ($\Delta$ \onAllCuaDelta{} points, 95\% CI \onAllCuaCi; \onAllCuaDisc{} tasks better and worse), and on the \textsc{test} split alone \onTestCuaArm{} against \onTestCuaOrig{} ($\Delta$ \onTestCuaDelta, CI \onTestCuaCi). On the same \textsc{test} tasks the pod-served original scores \onTestOrigScore/120, \onTestNoiseServerDelta{} points above the Modal-served one (CI \onTestNoiseServerCi): the noise between two servings of the same weights is as large as the merge's effect. % source: data/local/online/results_test_gh_merge_cua_b050, results_test_gh_orig_modal_test, results_confirm_gh_*_modal_confirm (figs/online.py)
ExPO was also run on the \textsc{test} split without having passed the gate, a coordinator's decision we report as such: \onTestExpoArm{} against \onTestExpoOrig{} ($\Delta$ \onTestExpoDelta{} points, 95\% CI \onTestExpoCi; \onTestExpoDisc{} tasks better and worse). All online scores are the official OSWorld score, a mean over tasks with partial credit included. With few tasks scoring differently between two runs (for the merge on the reserve, \onReserveCuaDisc{} better and worse), the bootstrap bounds move by up to a few points with the resampling seed; we report seed 0 throughout, and an independent recomputation agrees on every point estimate. counting only full successes, the \textsc{test} difference for ExPO is \onTestExpoBinDelta{} points (95\% CI \onTestExpoBinCi). % source: data/local/online/test_compare.json (online-eval)
On the reserve, a paired difference smaller than about \onMde{} points would go undetected with 80\% power, so these comparisons exclude only large effects.
A cheaper paired design, branching from recorded episodes, was piloted: replaying a recorded episode up to its last step reproduces the screenshots in \brValidOk{} of \brValidN{} validation episodes, and from the final step of \brEpisodes{} recorded episodes (two continuations each, both arms served by Modal) the original succeeds in \brOrig{} continuations and the merge in \brCand, with one episode whose replay diverges from its first step. It shows no advantage for the merge and is too small to show anything else. % source: data/local/online/branch_pilot/REPORT.md and {orig,cand}/branches.jsonl (figs/online.py)
\placefig{tab-b4-heldout}
